% !TeX program = luatex
% !TEX encoding = UTF-8


\documentclass{wg21}

\title{Safer \tcode{atomic\_ref::address} (FR-030-310)}
\docnumber{P3936R1}
\audience{LWG}
\author{Corentin Jabot}{corentin.jabot@gmail.com}
\authortwo{Gonzalo Brito Gadeschi}{gonzalob@nvidia.com}

\usepackage{color, colortbl}
\begin{document}
\maketitle

\section{Abstract}

This paper proposes to resolve FR-030-310 by returning a \tcode{void*} from \tcode{atomic_ref::address}


\section{Revisions}

\subsection{R1}

\begin{itemize}
    \item Wording fixes
    \item rename \tcode{address_return_t} and make it exposition only per LEWG guidance
\end{itemize}

\section{Motivation}

\tcode{atomic_ref<T>::address} currently returns a \tcode{T*}. This function was added by \paper{P2835R7},
which cover the use cases for this function.
In short, we might want to know the address for hashing for contention-aware data structure, indexing an array atomically, etc.

However, it is generally unsafe to access the object at \tcode{address} directly, and so, the NB comment argues that \tcode{address} is maybe too convenient.
IE, it is easy to use incorrectly, and hard to use correctly.
A main concern is that any other \tcode{atomic_ref} over the same object would no longer guarantee atomic access.

These observation were made in the original paper.

\begin{quoteblock}
The name and return type should prevent accidental misuse that could result from accessing the object through the address while there are still live \tcode{atomic_ref} that reference it.
\end{quoteblock}

However, of the options explored at the time, none seemed compelling.

\begin{itemize}
\item \tcode{uintptr\_t} is not mandated by the standard and would not work during constant evaluation.
\item \tcode{void*} would not allow us to get a \tcode{T*} back during constant evaluation.
\end{itemize}

FR-030-310 proposes to mandate \tcode{uintptr\_t}. However, EWG had no appetite to do it in the C++26 time frame (over concernns that there
might exist architectures in which pointers can be larger than the largest integer type). And it would not resolve the concerns
that it would not be usable during constant evaluation.

But... as it turns out, \tcode{void*} is now a suitable option. Indeed cast from \tcode{void*} to \tcode{T*} are now supported
in constant evaluation. This feature was added by \paper{P2738R1} in C++26.

The only downsize of using \tcode{void*} is then a slightly more complicated return type. However, that return type is really just copying the CV qualifiers,
which is not exactly novel.

Therefore, we think using \tcode{void*}  as the return type of \tcode{address} addresses the NB comment in a satisfactory manner.
\begin{itemize}
\item It addresses the concerns over missuse and safety: the \tcode{void*} pointer can be compared, hashed, etc without issue, but
accessing the underlying object requires an explicit cast.
\item Accessing the underlying object is still possible, even during constant evaluation.
\end{itemize}

\subsection{What about functions?}

Pointers to function are objects so this also works during constant evaluation \href{https://godbolt.org/z/h3bWxvPKn}{[Compiler Explorer]}.

\section{Wording}

\ednote {Add the following paragraph at the start of [atomics.general]}

\begin{addedblock}
Let \tcode{\placeholdernc{COPYCV}(FROM, TO)} be an alias for type \tcode{TO}
with the addition of \tcode{FROM}'s top-level cv-qualifiers.
\end{addedblock}

\ednote{That wording is a verbatim copy from [meta], LWG might want to factorize it out}.

\rSec3[atomics.ref.generic.general]{General}

\ednote{Modify in the synopsys}


\begin{codeblock}
using value_type = remove_cv_t<T>;
@\added{using \placeholder{address-return-type} = \placeholder{COPYCV(T, void)*}; //expos}@
static constexpr size_t required_alignment = implementation-defined;

// ...

constexpr @\changed{T*}{\placeholder{address-return-type}}@ address() const noexcept;
\end{codeblock}

\rSec3[atomics.ref.ops]{Operations}


\indexlibrarymember{address}{atomic_ref<T>}%
\begin{itemdecl}

constexpr @\changed{T*}{\placeholder{address-return-type}}@ address() const noexcept;
\end{itemdecl}

\begin{itemdescr}
\pnum
\returns
\tcode{ptr}.
\end{itemdescr}

\rSec3[atomics.ref.int]{Specializations for integral types}
\ednote{Modify in the synopsys}
\begin{codeblock}
using value_type = @\placeholder{integral-type}@;
using difference_type = value_type;
@\added{using \placeholder{address-return-type} = \placeholder{COPYCV(\placeholder{integral-type}, void)*}; //\expos}@
static constexpr size_t required_alignment = @\impdefx{required alignment for \tcode{atomic_ref} type's operations}@;

static constexpr bool is_always_lock_free = @\impdefx{whether a given \tcode{atomic} type's operations are always lock free}@;

// ...

constexpr @\changed{\placeholder{integral-type}*}{\placeholder{address-return-type}}@ address() const noexcept;
\end{codeblock}


\rSec3[atomics.ref.float]{Specializations for floating-point types}
\ednote{Modify in the synopsys}

\begin{codeblock}
using value_type = @\placeholdernc{floating-point-type}@;
using difference_type = value_type;
@\added{using \placeholder{address-return-type} = \placeholder{COPYCV(\placeholder{floating-point-type}, void)*}; //\expos}@
static constexpr size_t required_alignment = @\impdefx{required alignment for \tcode{atomic_ref} type's operations}@;


static constexpr bool is_always_lock_free = @\impdefx{whether a given \tcode{atomic} type's operations are always lock free}@;

// ...

constexpr @\changed{\placeholder{floating-point-type}*}{\placeholder{address-return-type}}@ address() const noexcept;
\end{codeblock}

\rSec3[atomics.ref.pointer]{Partial specialization for pointers}
\ednote{Modify in the synopsys}

\begin{codeblock}
using value_type = remove_cv_t<@\placeholder{pointer-type}@>;
using difference_type = ptrdiff_t;
@\added{using \placeholder{address-return-type} = \placeholder{COPYCV(\placeholder{pointer-type}, void)*}; //\expos}@

static constexpr size_t required_alignment = @\impdefx{required alignment for \tcode{atomic_ref} type's operations}@;

// ...

constexpr @\changed{\placeholder{pointer-type}*}{\placeholder{address-return-type}}@ address() const noexcept;
\end{codeblock}


\section{Acknowledgments}

We would like to thank Hana Dusíková and Michael Hava for their help and feedback on this paper.

\bibliographystyle{plain}
\bibliography{wg21, extra}

\renewcommand{\section}[2]{}%

\begin{thebibliography}{9}

\bibitem[N5008]{N5008}
Thomas Köppe
\emph{Working Draft, Standard for Programming Language C++}\newline
\url{https://wg21.link/N5008}


\end{thebibliography}

\end{document}
